% Section 4.4, from lectures/L04/appendix/c_exercises.md.

\section{Exercises}
\label{c4:sec:exercises}\label{c4:app:c}

\begin{aside}
\textbf{How to work these.} These exercises reinforce the concepts from \secref{c4:sec:raw} and
\secref{c4:sec:smart}. Sets~1 to~3 build one program with a raw-pointer factory; Set~4 turns the
same program into one with a smart-pointer factory. Build and run it with the Makefile from
\secref{c1:sec:tips}, which already supports the \file{include} and \file{source} directories
below.

The solutions are in \secref{sol:sec:c4}. Try each exercise yourself before looking at them.
\end{aside}

Create the following directory structure:

\begin{console}
Makefile
include/
    app/
        logic/
            logic.hpp
    driver/
        factory/
            esp32s3.hpp
            interface.hpp
            stub.hpp
        serial/
            esp32s3.hpp
            interface.hpp
            stub.hpp
source/
    main.cpp
\end{console}

\note When a class overrides a method that is marked \code{[[nodiscard]]} in the interface, repeat
the attribute on the overriding method as well. The attribute is not inherited; see the structure
of an interface in \secref{c3:app:b}.

% ------------------------------------------------------------------------------------------
\exerciseset{Serial Interface and Drivers}
\label{c4:set:1}

\exercise{Serial interface}{Code}
\label{c4:ex:1.1}
In this exercise you will design an interface \code{driver::serial::Interface} representing a
generic serial driver.

\task{Tasks}
Design a class named \code{driver::serial::Interface} that represents a generic serial
communication driver.

All methods except the destructor shall be pure virtual (\code{= 0}).

\task{a) Destructor}
Add a destructor that is:
\begin{itemize}
  \item Virtual.
  \item Marked \code{noexcept}.
  \item Implemented using \code{= default}.
\end{itemize}

\task{b) Initialization status}
Add a pure virtual method \code{isInitialized()} that:
\begin{itemize}
  \item Indicates whether the driver has been initialized (\code{true/false}).
  \item Does not modify the object (\code{const}).
  \item Cannot throw exceptions.
  \item Generates a warning if the return value is discarded (\code{[[nodiscard]]}).
\end{itemize}

\task{c) Transmitting data}
Add a pure virtual method \code{write()} that:
\begin{itemize}
  \item Transmits one byte of data.
  \item Takes a parameter of type \code{std::uint8_t}.
  \item Does not return a value.
  \item Cannot throw exceptions.
\end{itemize}

\task{d) Receiving data}
Add a pure virtual method \code{read()} that:
\begin{itemize}
  \item Attempts to read one byte from the serial driver.
  \item Takes a reference to a variable where the byte will be stored.
  \item Returns \code{true} if a byte was received, otherwise \code{false}.
  \item Cannot throw exceptions.
  \item Lets the caller discard the return value; the caller does not need to check whether a byte
    was received.
\end{itemize}

\exercise{Stub serial driver}{Code}
\label{c4:ex:1.2}
In this exercise you will implement a stub driver \code{driver::serial::Stub}.

A stub driver simulates hardware behavior and is useful for testing without real hardware.

The class shall:
\begin{itemize}
  \item Inherit from \code{driver::serial::Interface}.
  \item Be marked \code{final}.
\end{itemize}

\task{a) Member variables}
Add three private member variables:
\begin{itemize}
  \item The first member variable shall:
  \begin{itemize}
    \item Store the most recently transmitted byte.
    \item Have the type \code{std::uint8_t}.
    \item Be named \code{myLastByte}.
  \end{itemize}
  \item The second member variable shall:
  \begin{itemize}
    \item Indicate whether the driver is initialized.
    \item Have the type \code{bool}.
    \item Be named \code{myInitialized}.
  \end{itemize}
  \item The third member variable shall:
  \begin{itemize}
    \item Indicate whether a byte is available to read.
    \item Have the type \code{bool}.
    \item Be named \code{myHasData}.
  \end{itemize}
\end{itemize}

\task{b) Constructor}
Add a default constructor that:
\begin{itemize}
  \item Sets the driver as initialized.
  \item Indicates that no data is available to read.
  \item Is marked \code{noexcept}.
\end{itemize}

\task{c) Destructor}
Add a destructor that is:
\begin{itemize}
  \item Public.
  \item Marked \code{noexcept}.
  \item Implemented using \code{= default}.
\end{itemize}

\task{d) Implement interface methods}
Implement all methods required by the interface:
\begin{itemize}
  \item The method \code{isInitialized()} shall return the value stored in \code{myInitialized}.
  \item The method \code{write()} shall:
  \begin{itemize}
    \item Do nothing if \code{myInitialized} is \code{false}.
    \item Otherwise store the transmitted byte in \code{myLastByte} and set \code{myHasData} to
      \code{true}.
  \end{itemize}
  \item The method \code{read()} shall:
  \begin{itemize}
    \item Return \code{false} if \code{myInitialized} is \code{false} or \code{myHasData} is
      \code{false}.
    \item Otherwise copy the stored byte (\code{myLastByte}) into the output argument, set
      \code{myHasData} to \code{false}, and return \code{true}.
  \end{itemize}
\end{itemize}

\task{e) Disable copy and move semantics}
Delete the following functions (in the public section of the class):
\begin{itemize}
  \item Copy constructor.
  \item Move constructor.
  \item Copy assignment operator.
  \item Move assignment operator.
\end{itemize}

\exercise{Serial driver for ESP32-S3}{Code}
\label{c4:ex:1.3}
In this exercise you will create a hardware-specific serial driver \code{driver::serial::Esp32s3}.

This class represents a serial driver for the ESP32-S3 microcontroller.

The class shall:
\begin{itemize}
  \item Inherit from \code{driver::serial::Interface}.
  \item Be marked \code{final}.
\end{itemize}

\task{a) Member variables}
Add two private member variables:
\begin{itemize}
  \item The first shall store the transmit pin number.
  \item The second shall store the receive pin number.
  \item Both shall have the type \code{std::uint8_t} and be marked \code{const}.
  \item They shall be named \code{myTxPin} and \code{myRxPin}.
\end{itemize}

\task{b) Constructor}
Add an explicit constructor that:
\begin{itemize}
  \item Takes a transmit pin number.
  \item Takes a receive pin number.
  \item Stores them in the corresponding member variables.
  \item Is marked \code{noexcept}.
\end{itemize}

\task{c) Destructor}
Add a destructor that is:
\begin{itemize}
  \item Public.
  \item Marked \code{noexcept}.
  \item Declared with \code{override}.
\end{itemize}

\task{d) Interface methods}
Declare and implement the required interface methods.

For this exercise, it is sufficient to use placeholder implementations:
\begin{itemize}
  \item \code{isInitialized()} may always return \code{true}.
  \item \code{write()} shall print the transmitted byte in hexadecimal format. For example,
    transmitting byte \code{0xFF} on TX pin \code{17} shall print
    \code{Transmitting byte 0xFF via TX pin 17!}.
  \item \code{read()} may always return \code{false}.
\end{itemize}

\task{e) Disable copy and move semantics}
Delete the following functions (in the public section of the class):
\begin{itemize}
  \item Default constructor.
  \item Copy constructor.
  \item Move constructor.
  \item Copy assignment operator.
  \item Move assignment operator.
\end{itemize}

% ------------------------------------------------------------------------------------------
\exerciseset{Factory Interface with Raw Pointers}
\label{c4:set:2}

\exercise{Factory interface}{Code}
\label{c4:ex:2.1}
In this exercise you will design a factory interface \code{driver::factory::Interface}.

The purpose of this factory is to create serial drivers without exposing the concrete driver type
to the application logic.

\task{Tasks}
Design a class named \code{driver::factory::Interface}.

\task{a) Destructor}
Add a destructor that is:
\begin{itemize}
  \item Virtual.
  \item Marked \code{noexcept}.
  \item Implemented using \code{= default}.
\end{itemize}

\task{b) Factory method}
Add a pure virtual method \code{serial()} that:
\begin{itemize}
  \item Creates a serial driver.
  \item Takes a transmit pin number of type \code{std::uint8_t}.
  \item Takes a receive pin number of type \code{std::uint8_t}.
  \item Returns a pointer to \code{driver::serial::Interface}.
  \item Is marked \code{noexcept}.
  \item Generates a warning if the return value is discarded (\code{[[nodiscard]]}).
\end{itemize}

\exercise{ESP32-S3 factory}{Code}
\label{c4:ex:2.2}
In this exercise you will implement \code{driver::factory::Esp32s3}.

This factory shall create real drivers for ESP32-S3.

The class shall:
\begin{itemize}
  \item Inherit from \code{driver::factory::Interface}.
  \item Be marked \code{final}.
\end{itemize}

\task{a) Constructor and destructor}
Add:
\begin{itemize}
  \item A default constructor marked \code{noexcept}.
  \item A destructor marked \code{noexcept} and \code{override}.
\end{itemize}

\task{b) Factory method}
Implement the method \code{serial()} so that it:
\begin{itemize}
  \item Takes a transmit pin number and a receive pin number as input.
  \item Dynamically allocates a \code{driver::serial::Esp32s3}.
  \item Returns the object as a pointer to \code{driver::serial::Interface}.
\end{itemize}
Use the operator \code{new}.

\task{c) Disable copy and move semantics}
Delete the following functions (in the public section of the class):
\begin{itemize}
  \item Copy constructor.
  \item Move constructor.
  \item Copy assignment operator.
  \item Move assignment operator.
\end{itemize}

\exercise{Stub factory}{Code}
\label{c4:ex:2.3}
In this exercise you will implement \code{driver::factory::Stub}.

This factory shall create stub serial drivers.

The class shall:
\begin{itemize}
  \item Inherit from \code{driver::factory::Interface}.
  \item Be marked \code{final}.
\end{itemize}

\task{a) Constructor and destructor}
Add:
\begin{itemize}
  \item A default constructor marked \code{noexcept}.
  \item A destructor marked \code{noexcept} and \code{override}.
\end{itemize}

\task{b) Factory method}
Implement the method \code{serial()} so that it:
\begin{itemize}
  \item Takes a transmit pin number and a receive pin number as input.
  \item Ignores both pin numbers.
  \item Dynamically allocates a \code{driver::serial::Stub}.
  \item Returns the object as a pointer to \code{driver::serial::Interface}.
\end{itemize}
Use the operator \code{new}.

\task{c) Disable copy and move semantics}
Delete the following functions (in the public section of the class):
\begin{itemize}
  \item Copy constructor.
  \item Move constructor.
  \item Copy assignment operator.
  \item Move assignment operator.
\end{itemize}

% ------------------------------------------------------------------------------------------
\exerciseset{Application Logic with Raw Pointers}
\label{c4:set:3}

\exercise{Logic class using a factory}{Code}
\label{c4:ex:3.1}
In this exercise you will implement a logic class \code{app::logic::Logic} that uses a factory to
create its serial driver.

The application logic shall:
\begin{itemize}
  \item Own one serial driver.
  \item Periodically transmit a message.
  \item Attempt to read one received byte.
\end{itemize}

\task{a) Member variable}
Add a private member variable \code{mySerial} of type \code{driver::serial::Interface*}.

\task{b) Constructor}
Add a constructor that:
\begin{itemize}
  \item Takes a reference to \code{driver::factory::Interface}.
  \item Takes a transmit pin number of type \code{std::uint8_t}.
  \item Takes a receive pin number of type \code{std::uint8_t}.
  \item Is marked \code{noexcept}.
\end{itemize}
Inside the constructor, use the factory to create the serial driver.

\task{c) Destructor}
Add a destructor that:
\begin{itemize}
  \item Is marked \code{noexcept}.
  \item Deletes \code{mySerial}.
\end{itemize}

\task{d) \code{run()} method}
Add a method \code{run()} that:
\begin{itemize}
  \item Is marked \code{noexcept}.
  \item Runs continuously.
  \item Sends an incrementing byte value (\code{0}--\code{255}). After reaching \code{255}, the value
    wraps around to \code{0}.
  \item Attempts to read one byte into a local variable and prints it if received.
  \item Repeats forever.
  \item Delays execution by \code{100 ms} at the end of each loop iteration:
\end{itemize}

\begin{cppcode}
#include <chrono>
#include <thread>

constexpr std::uint8_t sleep_ms{100U};

std::this_thread::sleep_for(std::chrono::milliseconds(sleep_ms));
\end{cppcode}

\noindent Use local variables of type \code{std::uint8_t} to store the transmitted and received
bytes:

\begin{cppcode}
std::uint8_t txByte{}, rxByte{};
\end{cppcode}

\task{e) Disable copy and move semantics}
Delete the following functions (in the public section of the class):
\begin{itemize}
  \item Default constructor.
  \item Copy constructor.
  \item Move constructor.
  \item Copy assignment operator.
  \item Move assignment operator.
\end{itemize}

\exercise{Using the raw-pointer factory in \code{main}}{Code}
\label{c4:ex:3.2}
In this exercise you will test your application using both factories.

\task{a) ESP32-S3 factory}
In \file{main.cpp}:
\begin{itemize}
  \item Create an instance of \code{driver::factory::Esp32s3}.
  \item Create an instance of \code{app::logic::Logic} using that factory.
  \item Pass suitable pin numbers, for example \code{17U} and \code{18U}.
  \item Call \code{run()}.
\end{itemize}

\task{b) Stub factory}
Modify \file{main.cpp} so that \code{driver::factory::Stub} is used instead of
\code{driver::factory::Esp32s3}.

\task{c) Reflection}
Answer the following questions:
\begin{itemize}
  \item What code in \code{Logic} had to change when switching from the ESP32-S3 factory to the
    stub factory?
  \item What code in \code{main()} had to change?
\end{itemize}

% ------------------------------------------------------------------------------------------
\exerciseset{Factory Interface with Smart Pointers}
\label{c4:set:4}

\exercise{Updating the factory interface}{Code}
\label{c4:ex:4.1}
In this exercise you will modernize the factory interface by replacing raw pointers with
\code{std::unique_ptr}.

\task{Tasks}
Update \code{driver::factory::Interface} so that the method \code{serial()}:
\begin{itemize}
  \item Returns \code{std::unique_ptr<driver::serial::Interface>}.
  \item Still takes a transmit pin number of type \code{std::uint8_t}.
  \item Still takes a receive pin number of type \code{std::uint8_t}.
  \item Is marked \code{noexcept}.
  \item Is still marked \code{[[nodiscard]]}: discarding the returned \code{std::unique_ptr} would
    destroy the newly created object immediately.
\end{itemize}
Also include the header \code{<memory>} where needed.

\exercise{Updating the concrete factories}{Code}
\label{c4:ex:4.2}
In this exercise you will update both concrete factories to use smart pointers.

\task{a) ESP32-S3 factory}
Modify \code{driver::factory::Esp32s3} so that:
\begin{itemize}
  \item The method \code{serial()} returns \code{std::unique_ptr<driver::serial::Interface>}.
  \item The object is created using
    \code{std::make_unique<driver::serial::Esp32s3>(txPin, rxPin)}.
\end{itemize}

\task{b) Stub factory}
Modify \code{driver::factory::Stub} so that:
\begin{itemize}
  \item The method \code{serial()} returns \code{std::unique_ptr<driver::serial::Interface>}.
  \item The object is created using \code{std::make_unique<driver::serial::Stub>()}.
  \item The pin numbers are ignored.
\end{itemize}

\exercise{Updating the logic class}{Code}
\label{c4:ex:4.3}
In this exercise you will update \code{app::logic::Logic} to use smart pointers.

\task{a) Member variable}
Replace the raw pointer member \code{mySerial} with a \code{std::unique_ptr<driver::serial::Interface>}.

\task{b) Constructor}
Keep the constructor parameters the same as before, but initialize the member variable using the
factory's smart-pointer return value.

\task{c) Destructor}
Remove the destructor entirely.

\task{d) \code{run()} method}
Keep the behavior the same as before:
\begin{itemize}
  \item Send an incrementing byte value (\code{0}--\code{255}), wrapping around to \code{0} after
    \code{255}.
  \item Attempt to read one byte and print it if received.
  \item Delay execution by \code{100 ms} at the end of each loop iteration.
  \item Repeat forever.
\end{itemize}
Remember that member access through a smart pointer still uses the arrow operator \code{->}.
